% Ordinary amsthm statements: names, independent counters, resets and links.
\documentclass[oneside,openany]{bookbangla}
\setbanglamainfont{Kalpurush}
\begin{document}
\mainmatter
\chapter{উপপাদ্য}
\begin{theorem}[পায়রার খোপ নীতি]\label{thm:named}
তিনটি বস্তু দুটি খোপে রাখলে কোনো খোপে অন্তত দুটি বস্তু থাকে।
\end{theorem}
\begin{proof}প্রতি খোপে একটি করে রাখলে মাত্র দুটি বস্তু রাখা যায়।\end{proof}
\begin{example}\label{example:separate}একটি উদাহরণ।\end{example}
\begin{exercise}\label{exercise:separate}একটি সমস্যা।\end{exercise}
\section{পরের পরিচ্ছেদ}
\begin{theorem}\label{thm:continued}আরেকটি উপপাদ্য।\end{theorem}
\chapter{পরের অধ্যায়}
\begin{theorem}\label{thm:reset}নতুন অধ্যায়ের উপপাদ্য।\end{theorem}
\appendix
\chapter{পরিশিষ্ট}
\begin{theorem}\label{thm:appendix}পরিশিষ্টের উপপাদ্য।\end{theorem}
\autoref{thm:named}, \autoref{thm:continued}, \autoref{thm:reset},
\autoref{thm:appendix}, \autoref{example:separate}, \autoref{exercise:separate}।
\end{document}
